\subsubsection{二回分の更新}\label{節：周期係数を行列で見直す}
  \sectionref*{節：周期的な係数}では,$2x+1$の後に$3x+2$を計算し,$6x+5$を得た.
  行列で書くと
  \[
    A=\begin{pmatrix}2&1\\0&1\end{pmatrix},\qquad
    B=\begin{pmatrix}3&2\\0&1\end{pmatrix}.
  \]
  先に$A$,次に$B$なので,二回分は
  \[
    BA=\begin{pmatrix}6&5\\0&1\end{pmatrix},\qquad
    BA\binom{x}{1}=\binom{6x+5}{1}
  \]
  となる.この例では$AB=BA$だが,一般には異なる.接続演習問4なら
  \[
    U=\begin{pmatrix}-1&0\\0&1\end{pmatrix},\qquad
    V=\begin{pmatrix}2&1\\0&1\end{pmatrix}
  \]
  に対して,
  \[
    VU=\begin{pmatrix}-2&1\\0&1\end{pmatrix},\qquad
    UV=\begin{pmatrix}-2&-1\\0&1\end{pmatrix}
  \]
  となる.更新を行う順序が,付加項の符号に現れた.
  一般の$a_{n+1}=p_na_n+q_n$でも,$\binom{a_{n+1}}1=
  \begin{pmatrix}p_n&q_n\\0&1\end{pmatrix}\binom{a_n}1$と表せる.
  一周期分は,先に行う\bookterm{matrix}{行列}を右に置いて掛け合わせる.途中の値は一回ずつの更新で求める.
  \bookproblemref{practice-connections}{4}{接続演習問4}で本文との対応を確かめよう.
